\documentclass[11pt]{article}
\newcommand{\DocShort}{Anexo A}
\usepackage{fontspec}
\setmainfont{TeX Gyre Pagella}
\setsansfont{TeX Gyre Heros}
\usepackage[spanish,es-noshorthands]{babel}
\decimalpoint
\usepackage{geometry}
\geometry{letterpaper, top=2.4cm, bottom=2.4cm, left=2.7cm, right=2.7cm}
\usepackage[expansion=false]{microtype}
\usepackage{parskip}
\setlength{\parskip}{6pt}
\usepackage{amsmath,amssymb}
\usepackage{booktabs}
\usepackage{array}
\usepackage{longtable}
\usepackage{caption}
\usepackage[dvipsnames,table]{xcolor}
\usepackage{enumitem}
\setlist[itemize]{topsep=2pt, itemsep=2pt, parsep=0pt}
\setlist[enumerate]{topsep=2pt, itemsep=2pt, parsep=0pt}
\usepackage{fancyhdr}
\usepackage[hidelinks,pdfencoding=auto]{hyperref}
\usepackage{titlesec}
\usepackage{tcolorbox}
\tcbuselibrary{breakable}
\usepackage{tabularx}
\usepackage{makecell}
\usepackage{ragged2e}
\usepackage{csquotes}
\usepackage[backend=biber,style=apa,sorting=nyt,sortcites=true,language=spanish]{biblatex}
\DeclareLanguageMapping{spanish}{spanish-apa}
\addbibresource{referencias_presentaciones.bib}

\definecolor{darkblue}{HTML}{174A63}
\definecolor{medblue}{HTML}{2E7896}
\definecolor{lightblue}{HTML}{DCEEF4}
\definecolor{teal}{HTML}{1B7F79}
\definecolor{lightteal}{HTML}{D7EFEC}
\definecolor{lightgrey}{HTML}{F5F7F8}
\definecolor{midgrey}{HTML}{5B6770}
\definecolor{warnyellow}{HTML}{FFF4CC}

\titleformat{\section}
{\large\bfseries\color{darkblue}}
{\thesection}{0.75em}{}[{\color{medblue}\titlerule[0.8pt]}]
\titleformat{\subsection}
{\normalsize\bfseries\color{darkblue}}
{\thesubsection}{0.75em}{}
\titleformat{\subsubsection}
{\normalsize\bfseries\color{medblue}}
{\thesubsubsection}{0.75em}{}

\pagestyle{fancy}
\fancyhf{}
\fancyhead[L]{\small\color{midgrey}ICV-342-T Hidrología}
\fancyhead[R]{\small\color{midgrey}\DocShort}
\fancyfoot[C]{\small\color{midgrey}\thepage}
\setlength{\headheight}{14pt}
\renewcommand{\headrulewidth}{0.4pt}
\renewcommand{\footrulewidth}{0pt}

\rowcolors{2}{lightgrey}{white}
\renewcommand{\arraystretch}{1.18}
\newcolumntype{Y}{>{\RaggedRight\arraybackslash}X}

\newtcolorbox{infobox}{
  breakable,
  colback=lightblue,
  colframe=medblue,
  boxrule=0.8pt,
  arc=3pt,
  left=8pt,
  right=8pt,
  top=7pt,
  bottom=7pt
}

\newtcolorbox{warnbox}{
  breakable,
  colback=warnyellow,
  colframe=orange!70!black,
  boxrule=0.8pt,
  arc=3pt,
  left=8pt,
  right=8pt,
  top=7pt,
  bottom=7pt
}

\newcommand{\doccover}[2]{
  \begin{center}
  {\small\color{midgrey}Pontificia Universidad Católica Madre y Maestra\\Escuela de Ingeniería Civil y Ambiental}

  \vspace{1.0cm}
  {\color{medblue}\rule{\linewidth}{2pt}}
  \vspace{0.3cm}

  {\Large\bfseries\color{darkblue}ICV-342-T Hidrología}\\[0.35em]
  {\LARGE\bfseries\color{darkblue}#1}

  \vspace{0.3cm}
  {\color{medblue}\rule{\linewidth}{2pt}}
  \vspace{0.9cm}
  \end{center}

  \begin{center}
  \begin{tabular}{ll}
  \toprule
  \textbf{Asignatura} & ICV-342-T Hidrología \\
  \textbf{Profesor} & Ricardo Rafael Hernández Moreira, PhD \\
  \textbf{Periodo académico} & Mayo--agosto de 2026 \\
  \textbf{Versión} & #2 \\
  \bottomrule
  \end{tabular}
  \end{center}

  \vspace{0.4cm}
}

\newcommand{\unitbox}[3]{
  \begin{tcolorbox}[
    breakable,
    colback=lightteal,
    colframe=teal,
    boxrule=0.8pt,
    arc=3pt,
    title=\textbf{#1}
  ]
  \textbf{Enfoque de la exposición.} #2

  \vspace{0.5\baselineskip}
  \textbf{Apoyo bibliográfico sugerido.} #3
  \end{tcolorbox}
}

\AtBeginDocument{
  \hypersetup{
    pdftitle={\DocShort},
    pdfauthor={Ricardo Rafael Hernández Moreira},
    pdfsubject={ICV-342-T Hidrología}
  }
}


\begin{document}
\doccover{Anexo A. Formulario breve de audiencia y revisión de pares}{26 de mayo de 2026}

\begin{infobox}
Este formulario es individual. Sirve para controlar asistencia, atención y comprensión mínima de la sesión. Puede ser compartido con el equipo expositor sin identificar al estudiante.
\end{infobox}

\section{Rúbrica de audiencia}
\begin{itemize}
\item La audiencia no es pasiva. Cada estudiante debe asistir, seguir la exposición y entregar este formulario al final de la sesión.
\item La parte de audiencia equivale a 1\% del curso, distribuida entre todas las sesiones de exposición.
\item El formulario debe llevar nombre y matrícula para que el profesor pueda registrar asistencia y calificación. Si se comparten observaciones con el equipo expositor, se hará sin identificar al estudiante que las escribió.
\item Si un estudiante falta a una sesión de exposiciones o no entrega el formulario correspondiente, obtiene cero en esa sesión para la parte de audiencia.
\end{itemize}

\begin{center}
\begin{tabular}{p{0.50\linewidth} p{0.34\linewidth}}
\toprule
\textbf{Criterio de audiencia} & \textbf{Se espera} \\
\midrule
Asistencia y entrega (40\%) & Presencia efectiva y entrega del formulario de la sesión. \\
Atención verificable (35\%) & Comentario, observación o pregunta conectada con el contenido expuesto. \\
Respuesta técnica breve (25\%) & Respuesta clara a una pregunta técnica corta sobre la exposición. \\
\bottomrule
\end{tabular}
\end{center}

\section{Formulario}
\begin{tabularx}{\linewidth}{p{0.28\linewidth} Y}
\toprule
\textbf{Campo} & \textbf{Completar} \\
\midrule
Nombre del estudiante & \\
Matrícula & \\
Fecha & \\
Unidad expuesta & \\
Equipo expositor & \\
\bottomrule
\end{tabularx}

\subsection*{1. Verificación rápida}
\begin{tabularx}{\linewidth}{Y p{0.20\linewidth}}
\toprule
\textbf{Ítem} & \textbf{Respuesta} \\
\midrule
Asistí a la sesión completa de exposición. & Sí / No \\
Tomé nota de al menos una idea o dato técnico relevante. & Sí / No \\
Formulé una pregunta, observación o comentario técnico, o puedo dejarlo por escrito abajo. & Sí / No \\
\bottomrule
\end{tabularx}

\subsection*{2. Idea técnica principal de la sesión}
\vspace{2.4cm}

\subsection*{3. Observación o pregunta técnica para el equipo expositor}
\vspace{2.8cm}

\subsection*{4. Pregunta técnica breve}
\textit{Redactada por el profesor o tomada del núcleo de la exposición. Debe poder responderse en 3 a 5 líneas.}

\vspace{3.0cm}

\subsection*{5. Valoración breve de la exposición}
\begin{tabularx}{\linewidth}{Y p{0.18\linewidth}}
\toprule
\textbf{Criterio} & \textbf{Nivel} \\
\midrule
La exposición fue clara y ordenada. & \makecell[l]{Alta /\\ Media /\\ Baja} \\
El equipo mostró dominio técnico razonable. & \makecell[l]{Alta /\\ Media /\\ Baja} \\
El material visual ayudó a comprender el tema. & \makecell[l]{Alta /\\ Media /\\ Baja} \\
\bottomrule
\end{tabularx}

\begin{warnbox}
La entrega de este formulario sin contenido técnico verificable puede ser registrada como participación insuficiente.
\end{warnbox}

\end{document}
